\section{Conclusion}
This project developed Minimax, Monte Carlo Tree Search (MCTS) and tabular Q-learning (TabQL) agents for tic-tac-toe and the complete two-player Azul game. Each baseline agent worked in tic-tac-toe, and every game in the performed head-to-head matches ended in a draw. These results indicate strong play, and are potentially indicative of perfect play for MCTS and TabQL.\\

Moving to Azul required adaptations for random factory refills and a much larger number of possible moves. In the Azul trials described in this report, Minimax and MCTS produced encouraging results. Enhancements to agents generally reduced search time or improved net score. TabQL was also implemented for Azul, but under the training and evaluation conditions used, it did not consistently beat the random agent and did not beat Minimax. This result is likely due to the large state and action space that Azul has.\\

The shared environment interface has been designed so that further turn-based, two-player stochastic games could be added by implementing the required game operations and providing suitable game-specific components. Time constraints prevented implementations using Deep Q-Networks (DQN) and Proximal Policy Optimisation (PPO), leaving the reinforcement-learning comparison limited to TabQL. Future work could investigate those methods, explore an AlphaGo-inspired\cite{alphagozero} combination of MCTS and reinforcement learning, and test the interface with further games. Connect Four could provide an intermediate step between tic-tac-toe and Azul, while chess could offer a separate, more demanding test.